\chapter{\ploren{Podsumowanie}{Conclusions}}

W ostatnim rozdziale zamieszczamy podsumowanie informacji zawartych w głównej części pracy, w tym m.in. zalety i ograniczenia proponowanych metod i rozwiązań. Jeszcze raz podkreślić należy własny wkład autora, zaś w przypadku pracy magisterskiej – wskazać na element nowatorstwa. W podsumowaniu nie wracamy już do zagadnień teoretycznych, zwykle też nie umieszczamy rysunków, tabel, itp.

Dyskusja nad dalszym rozwojem pracy. Wnioski. Omówienie wyników. Co zrobiono w pracy i jakie uzyskano wyniki? Czy i w jakim zakresie praca stanowi nowe ujęcie problemu? Sposób wykorzystania pracy (publikacja, udostępnienie instytucjom, materiał źródłowy dla studentów). Co uważa autor za własne osiągnięcia? Perspektywy wdrożenia, możliwości i propozycje kontynuacji. Ten rozdział zajmie około 2-3 stron.
